% ---------------------------------------------------------------------------
% formose_chemistry.tex --- the prescribed formose routes as one convergent
% scheme: three glycolaldehyde molecules, two enolizations (r1), the tetrose
% aldol (r3), and the hexose aldol (r7).
%
% Source fragment: loaded by SI/si.tex.
%
% Skeletal structures, hydrogens on carbon implicit.  No atom mapping and no
% stereochemical outcome is asserted: the rule-generated network does not
% assign stereocentres.  Reactions are drawn as Petri-net bars, as in the main
% text, so each glycolaldehyde molecule is a separate input; the shaded block
% is the T4 route (r1 + r3), the whole scheme the T6 route (2 r1 + r3 + r7).
% Black and white, after the shared result-figure style.
% ---------------------------------------------------------------------------
% Figure styles are supplied by the manuscript preamble.

\begingroup
\setchemfig{atom sep=17pt,bond style={line width=0.5pt},bond offset=1.1pt}
\renewcommand*\printatom[1]{\ensuremath{\mathsf{#1}}}
\newcommand{\mfont}{\fontsize{6.5}{7.5}\selectfont\color{black}}
\newcommand{\glycolald}{\chemfig{HO-[:30]-[:-30](-[:30]H)=[:-90]O}}
\newcommand{\enediol}{\chemfig{HO-[:30]=[:-30]-[:30]OH}}
\newcommand{\tetrose}{\chemfig{HO-[:30]-[:-30](-[:-90]OH)-[:30](-[:90]OH)-[:-30](-[:30]H)=[:-90]O}}
\newcommand{\hexose}{\chemfig{HO-[:30]-[:-30](-[:-90]OH)-[:30](-[:90]OH)-[:-30](-[:-90]OH)-[:30](-[:90]OH)-[:-30](-[:30]H)=[:-90]O}}

\begin{tikzpicture}[syn]
\def\W{16.40}
\path[use as bounding box] (0,0) rectangle (\W,-5.25);

\tikzset{
  mol/.style={anchor=center,inner sep=1pt},
  sname/.style={font=\scSmall,text=black,anchor=north,align=center},
  feed/.style={-{Stealth[round,length=1.9mm,width=1.5mm]},line width=0.6pt,
               draw=black},
  rlab/.style={font=\scSmall,text=black,anchor=south,inner sep=1.5pt},
  note/.style={font=\scSmall,text=black,align=left,anchor=north west},
}
% reaction bar: \rbar{name}{x}{y}{label}
\newcommand{\rbar}[4]{%
  \coordinate (#1) at (#2,#3);
  \draw[line width=2.4pt] (#2,#3-0.30) -- (#2,#3+0.30);
  \node[rlab] at (#2,#3+0.30) {#4};
  \coordinate (#1w) at (#2-0.06,#3);
  \coordinate (#1e) at (#2+0.06,#3);}

% Main chain on one line: glycolaldehyde -> r3 -> T4 -> r7 -> T6.  Each aldol
% takes its enediol from a separate enolization, fed from above (r3) or below
% (r7), so the three glycolaldehyde inputs are drawn as three molecules.
\def\ym{-2.35}   % main chain
\def\yt{-0.85}   % first enolization
\def\yl{-4.30}   % second enolization

% --- T4 route (r1 + r3): shaded block --------------------------------------
\fill[black!7] (0.15,-0.12) rectangle (9.85,-3.62);
\node[note] at (0.25,-2.55) {$T_4$ route: $r_1+r_3$\\2 firings\\2 glycolaldehyde};

% --- first enolization (top) and tetrose aldol (main chain) ----------------
\node[mol] (ga) at (1.55,\yt) {\mfont\glycolald};
\rbar{ra}{3.05}{\yt}{$r_1$}
\node[mol] (ea) at (4.40,\yt) {\mfont\enediol};
\draw[feed] (ga.east) -- (raw);
\draw[feed] (rae) -- (ea.west);

\node[mol] (gb) at (4.40,\ym) {\mfont\glycolald};
\rbar{rb}{6.10}{\ym}{$r_3$}
\draw[feed] (ea.east) to[out=0,in=150] (rbw);
\draw[feed] (gb.east) -- (rbw);
\node[mol] (t4) at (8.10,\ym) {\mfont\tetrose};
\draw[feed] (rbe) -- (t4.west);

% --- second enolization (bottom) and hexose aldol (main chain) -------------
\node[mol] (gc) at (5.20,\yl) {\mfont\glycolald};
\rbar{rc}{6.70}{\yl}{$r_1$}
\node[mol] (ec) at (8.10,\yl) {\mfont\enediol};
\draw[feed] (gc.east) -- (rcw);
\draw[feed] (rce) -- (ec.west);

\rbar{rd}{10.55}{\ym}{$r_7$}
\draw[feed] (t4.east) -- (rdw);
\draw[feed] (ec.east) to[out=0,in=-150] (rdw);
\node[mol] (t6) at (13.40,\ym) {\mfont\hexose};
\draw[feed] (rde) -- (t6.west);

% --- names under each structure ---------------------------------------------
\foreach \n/\t in {ga/glycolaldehyde 1, gb/glycolaldehyde 2,
                   gc/glycolaldehyde 3, ea/enediol, ec/enediol}
  \node[sname] at ([yshift=-0.5mm]\n.south) {\t};
\node[sname] at ([yshift=-0.5mm]t4.south) {$T_4$, aldotetrose (C$_4$)};
\node[sname] at ([yshift=-0.5mm]t6.south) {$T_6$, aldohexose (C$_6$)};

\node[note,anchor=north east,align=right] at (16.30,-3.75)
  {$T_6$ route: $2r_1+r_3+r_7$\\4 firings\\3 glycolaldehyde};
\end{tikzpicture}
\endgroup
